\documentclass[11pt]{article}
\usepackage[margin=1in]{geometry}
\usepackage{amsmath,amssymb,amsthm}
\usepackage{graphicx}
\usepackage[hidelinks]{hyperref}
\emergencystretch=3em

\newtheorem{theorem}{Theorem}
\newtheorem{lemma}{Lemma}
\newtheorem{remark}{Remark}

\newcommand{\C}{\mathbb C}
\newcommand{\Z}{\mathbb Z}
\newcommand{\E}{\mathbb E}
\newcommand{\Tr}{\operatorname{Tr}}

\title{Physical-Space Matrix Hypercontractivity over Large Alphabets for Target Exponents at Least Two}
\author{Run fa\_banach\_001, lane 19}
\date{July 18, 2026}

\begin{document}
\maketitle

\section*{Source Question}

Arunachalam and Doriguello ask in Future work item 4 of
arXiv:2109.02600 whether one can prove a general $(q,p)$-hypercontractive
statement for matrices, first for matrix-valued functions over $\{0,1\}^n$,
and then for functions over $\Z_r^n$. The source passage is shown in
Figure~\ref{fig:open}.

\begin{figure}[h]
\centering
\includegraphics[width=\textwidth]{figures/open_problem_crop.png}
\caption{Future work item 4 in arXiv:2109.02600, PDF page 9.}
\label{fig:open}
\end{figure}

This packet proves a physical-space theorem in the range
$1\leq p\leq 2\leq q<\infty$. It is therefore stronger than the earlier
coefficient-space packet for this target, but still does not settle the
intermediate range $p<q<2$.

\section*{Notation}

Let $G=\Z_r^n$ with normalized counting measure and characters
$\chi_S(x)=\omega_r^{S\cdot x}$, where $\omega_r=e^{2\pi i/r}$. For
$f:G\to\C^{m\times m}$ write
\[
  \widehat f(S)=\E_{x\in G} f(x)\overline{\chi_S(x)}, \qquad
  T_\rho f=\sum_{S\in G}\rho^{|S|}\widehat f(S)\chi_S .
\]
Here $|S|=|\{j:S_j\neq0\}|$. The mixed norm is
\[
  \|f\|_{L_s(S_s)}
  =
  \left(\E_{x\in G}\|f(x)\|_s^s\right)^{1/s},
\]
where $\|\cdot\|_s$ is the unnormalized Schatten $s$-norm. For
$1\leq s\leq 2$ set
\[
  \zeta_r(s)=
  \frac{(s-1)\bigl(1-(s-1)^{r-1}\bigr)}{(r-1)(2-s)}
\]
with the continuous values $\zeta_r(1)=0$ and $\zeta_r(2)=1$. For
$2\leq q<\infty$ set
\[
  \eta_r(q)=
  \frac{(q-1)\bigl((q-1)^{r-1}-1\bigr)}{(r-1)(q-2)}
\]
with $\eta_r(2)=1$.

\section*{Proof Intuition}

The paper proves an $r$-point Ball--Carlen--Lieb convexity inequality:
for $1\leq p\leq2$, the $L_p(S_p)$ norm of one $\Z_r$ coordinate controls
a weighted Hilbertian sum of its matrix Fourier coefficients, with weight
$\zeta_r(p)$ on all nonzero frequencies.

For $q\geq2$ the Ball--Carlen--Lieb smoothness inequality gives the reverse
estimate with constant $\eta_r(q)$. Applying this smoothness estimate to the
noisy coefficients, using $\|A\|_q\leq\|A\|_p$, and then applying the source
convexity estimate yields the one-coordinate $L_p(S_p)\to L_q(S_q)$
contraction whenever $\eta_r(q)\rho^2\leq\zeta_r(p)$. The product result is
obtained by treating all other coordinates as a finite direct sum of matrix
algebras.

\section*{Auxiliary Inequalities}

\begin{lemma}[Source $r$-point convexity]
\label{lem:convexity}
Let $1\leq p\leq2$ and let $A_0,\ldots,A_{r-1}$ be matrices, or elements of
a finite direct sum of matrix algebras. Then
\[
  \left(
  \frac1r\sum_{x=0}^{r-1}
  \left\|\sum_{k=0}^{r-1}\omega_r^{kx}A_k\right\|_p^p
  \right)^{2/p}
  \geq
  \|A_0\|_p^2+\zeta_r(p)\sum_{k=1}^{r-1}\|A_k\|_p^2 .
\]
\end{lemma}

\begin{proof}
For matrices this is exactly the $r$-point Ball--Carlen--Lieb inequality
proved in Section 3.1 of arXiv:2109.02600. The same statement for finite
direct sums follows by applying the matrix statement to the corresponding
block-diagonal embedding; rescaling the trace only rescales both sides by the
same homogeneous factor.
\end{proof}

\begin{lemma}[$r$-point smoothness]
\label{lem:smoothness}
Let $2\leq q<\infty$ and let $A_0,\ldots,A_{r-1}$ be matrices, or elements of
a finite direct sum of matrix algebras. Then
\[
  \left(
  \frac1r\sum_{x=0}^{r-1}
  \left\|\sum_{k=0}^{r-1}\omega_r^{kx}A_k\right\|_q^q
  \right)^{2/q}
  \leq
  \|A_0\|_q^2+\eta_r(q)\sum_{k=1}^{r-1}\|A_k\|_q^2 .
\]
\end{lemma}

\begin{proof}
We mimic the block-diagonal proof of the source convexity inequality, using
the Ball--Carlen--Lieb smoothness inequality in place of convexity. That
smoothness inequality says that if the same trace-orthogonality condition used
in the source proof holds for $Z$ and $W$, then
\[
  \|Z+W\|_q^2\leq \|Z\|_q^2+(q-1)\|W\|_q^2 .
\]
For $j=0,\ldots,r-1$, define the $r$-block matrices
\[
  Z_j=\operatorname{diag}\bigl(A_j,\omega_r^jA_j,\ldots,
  \omega_r^{(r-1)j}A_j\bigr).
\]
As in the source proof, the root-of-unity cancellation gives the required
trace-orthogonality between $Z_j$ and $Z_k$ whenever $j\neq k$, and also
between each selected $Z_j$ and the remaining tail.

For any permutation $(k_1,\ldots,k_{r-1})$ of $\{1,\ldots,r-1\}$, repeated
application of smoothness gives
\[
  \left\|\sum_{j=0}^{r-1}Z_j\right\|_q^2
  \leq
  \|Z_0\|_q^2+\sum_{\ell=1}^{r-1}(q-1)^\ell\|Z_{k_\ell}\|_q^2 .
\]
Averaging over all permutations yields
\[
  \left\|\sum_{j=0}^{r-1}Z_j\right\|_q^2
  \leq
  \|Z_0\|_q^2+\eta_r(q)\sum_{j=1}^{r-1}\|Z_j\|_q^2 .
\]
Finally,
\[
  \left\|\sum_j Z_j\right\|_q^q
  =
  \sum_{x=0}^{r-1}
  \left\|\sum_{j=0}^{r-1}\omega_r^{jx}A_j\right\|_q^q,
  \qquad
  \|Z_j\|_q^2=r^{2/q}\|A_j\|_q^2.
\]
Dividing by $r^{2/q}$ gives the displayed inequality. The finite direct-sum
case follows by the same block-diagonal embedding.
\end{proof}

\section*{Main Theorem}

\begin{theorem}[Physical-space matrix hypercontractivity for $q\geq2$]
\label{thm:main}
Let $r\geq2$, $n,m\geq1$, and $1\leq p\leq2\leq q<\infty$. If
\[
  0\leq \rho\leq \left(\frac{\zeta_r(p)}{\eta_r(q)}\right)^{1/2},
\]
then every $f:\Z_r^n\to\C^{m\times m}$ satisfies
\[
  \|T_\rho f\|_{L_q(S_q)}\leq \|f\|_{L_p(S_p)} .
\]
For $r=2$, this is the classical threshold
$\rho\leq\sqrt{(p-1)/(q-1)}$.
\end{theorem}

\begin{proof}
First consider one coordinate, so $f:\Z_r\to\mathcal A$ where $\mathcal A$ is
a finite direct sum of matrix algebras. Write
$f(x)=\sum_{k=0}^{r-1}\omega_r^{kx}A_k$. By Lemma~\ref{lem:smoothness},
\[
  \|T_\rho f\|_{L_q(S_q)}^2
  \leq
  \|A_0\|_q^2+\eta_r(q)\rho^2\sum_{k=1}^{r-1}\|A_k\|_q^2.
\]
Since $q\geq p$, Schatten monotonicity gives $\|A_k\|_q\leq\|A_k\|_p$.
The assumed bound on $\rho$ therefore implies
\[
  \|T_\rho f\|_{L_q(S_q)}^2
  \leq
  \|A_0\|_p^2+\zeta_r(p)\sum_{k=1}^{r-1}\|A_k\|_p^2.
\]
Lemma~\ref{lem:convexity} bounds the last expression by
$\|f\|_{L_p(S_p)}^2$. This proves the one-coordinate estimate with values in
any finite direct sum of matrix algebras.

For $G=\Z_r^n$, write $T_\rho=T_\rho^{(1)}\cdots T_\rho^{(n)}$, where
$T_\rho^{(j)}$ acts on the $j$th coordinate. The single-coordinate noise
kernel is the probability measure assigning mass
$(1+(r-1)\rho)/r$ to $0$ and $(1-\rho)/r$ to each nonzero element, so
$T_\rho^{(j)}$ is an $L_p(S_p)$ contraction for $0\leq\rho\leq1$.

Apply the one-coordinate estimate to the last coordinate, viewing the first
$n-1$ coordinates and the matrix variable as one finite direct sum of matrix
algebras. Since the remaining coordinate noises commute with the last one and
are $L_p(S_p)$ contractions,
\[
\begin{aligned}
  \|T_\rho f\|_{L_q(S_q)}
  &=
  \left\|T_\rho^{(n)}T_\rho^{(1)}\cdots T_\rho^{(n-1)}f\right\|_{L_q(S_q)}\\
  &\leq
  \left\|T_\rho^{(1)}\cdots T_\rho^{(n-1)}f\right\|_{L_p(S_p)}
  \leq \|f\|_{L_p(S_p)} .
\end{aligned}
\]
This proves the theorem.
\end{proof}

\section*{Scope and Remaining Range}

The theorem is a physical-space matrix-valued hypercontractivity statement
over $\Z_r^n$ for every $1\leq p\leq2\leq q<\infty$. It therefore answers the
large-target-exponent part of the source question with an explicit
dimension-free constant. It also recovers the optimal scalar threshold in the
Boolean case $r=2$ for this range.

The theorem does not settle $p<q<2$. The simple smoothness/convexity sandwich
used above is intrinsically pinned at the Hilbert exponent $2$. Extending it
below $2$ appears to require sharper Schatten Hanner-type information; this is
exactly the sort of terrain the source paper flags when discussing why matrix
Hanner inequalities are only known in restricted ranges. Small random $2\times2$
tests did not reveal a counterexample, so this packet records the result as a
partial theorem rather than a disproof of the remaining range.

\section*{References}

\begin{enumerate}
\item S. Arunachalam and J. F. Doriguello,
\emph{Matrix hypercontractivity, streaming algorithms and LDCs: the large
alphabet case}, arXiv:2109.02600; ACM Transactions on Computation Theory
16(4), Article 21, 2024.
\item K. Ball, E. A. Carlen, and E. H. Lieb,
\emph{Sharp uniform convexity and smoothness inequalities for trace norms},
Inventiones Mathematicae 115 (1994), 463--482.
\end{enumerate}

\end{document}
